\documentclass[11pt,a4paper]{article}
\usepackage[margin=23mm]{geometry}
\usepackage{amsmath,amssymb,amsthm}
\setlength{\emergencystretch}{2em}
\newtheorem{lemma}{Lemma}
\newtheorem{theorem}[lemma]{Theorem}
\newcommand{\gen}[1]{\langle#1\rangle}
\newcommand{\Alt}{\operatorname{Alt}}
\title{The terminal singleton supplies an opposite head}
\author{Human Compression project: L044 continuation}
\date{}
\begin{document}\maketitle
Retain the complete critical six-point-head configuration and its
acquired proofs. Thus $G\le\Alt(\Omega)$ is transitive with connected
string C-marking, all maximal omissions are intransitive, and
$n=2r\ge18$. The first three marks are an actual supported
extendable six-point head on $D$, with exit vertex $v$ and
\[
          Q=G_{\ge3}=\Alt(W),\qquad
          W=\Omega\setminus(D\setminus\{v\}).
\]
Assume the terminal omission is the actual natural group
\[
          P=G_{\le r-2}=\Alt(\Omega\setminus\{z\}),       \tag{1}
\]
and every colour $3,\ldots,r-2$ has two full ambient-eligible
edges. This last hypothesis is explicit; it holds when the
terminal colour is the first unique colour after the head,
after the acquired exclusions of colours $3,4$.
The induction that produced the critical setting remains inherited.
No new rank induction is used below.
Put
\[
 \eta=\rho_{r-1},\quad \sigma=\rho_{r-2},\quad
 H=G_{\le r-3},\quad K=G_{\le r-4},\quad a=z^\eta.
\]
Transitivity gives $a\ne z$. Since $\eta$ centralizes $H$
and $H$ fixes $z$, it also fixes $a$. Thus
$\eta=(z\ a)\epsilon$, where $\epsilon$ is an odd involution
fixing $z,a$ and centralizing $H$. The terminal mark centralizes
the actual head, so fixes $D$ pointwise. Hence $z,a$ and
the support of $\epsilon$ lie outside $D$.

\begin{lemma}[The actual last mark is a double transposition]
There are distinct points $u,v'$ outside $D\cup\{z,a\}$ such that
\[
                       \eta=(z\ a)(u\ v'),
\]
and $K$ fixes $u,v'$ pointwise. The group $H$ preserves their pair.
\end{lemma}
\begin{proof}
Apply the complete odd-singleton theorem to $Q$ on $W$, with
the entire marked string reversed. Its rank is $p=r-3\ge6$
and its degree is $n-5=2p+1$. The actual first omission is
\[
               Q\cap P=\Alt(W\setminus\{z\}),
\]
by (1) and full ambient IP. Restriction to $W$ is faithful,
so IP, nonidentity marks, rank and connectedness are retained.
Every maximal omission is intransitive: a transitive omission
on $W$, together with the left head meeting $W$ only at $v$,
would make its corresponding ambient omission transitive.
The required reversed colours at least two are precisely the
original colours $3,\ldots,r-3$. Their ambient-eligible edges
lie in $W$ and remain eligible for the smaller $Q$-omissions.
This verifies every hypothesis of that theorem.

Write $H_Q=\gen{\rho_3,\ldots,\rho_{r-3}}$ and
$K_Q=\gen{\rho_3,\ldots,\rho_{r-4}}$.
The theorem gives either at least three $H_Q$-fixed points
in $W$, or the literal residual
$\eta=(z\ a)(u\ v')$, whose two-point orbit
$E=\{u,v'\}$ is fixed pointwise by $K_Q$.
The latter fixation is proved in the theorem's complete
second proof, including its localized contradiction when
$K_Q$ moves that pair. Since $E$ is disjoint from $D$,
the first three marks fix it as well, so $K$ fixes it.
This establishes the lemma in the literal residual.

Suppose instead that $H_Q$ has at least three fixed points.
None is $v$, because its literal first original mark
$\rho_3$ moves $v$. The left head fixes $W\setminus\{v\}$,
so these are also fixed by $H$. Put $U=\operatorname{Supp}(H)$,
$F=\Omega\setminus U$ and $d=r-2$, so $|F|\ge3$.
The pure-prefix theorem gives a head orbit $A$ of $H$
carrying pure $\Alt(A)$. Its exact intersection with $Q$
removes five head points and leaves all exterior orbits unchanged.
If $q$ counts the nontrivial $H$-orbits, the paired interval
$H\cap Q$ of rank $d-3$ gives
\[
          2(d-3)\le |U|-5-q,\qquad |U|\ge2d-1+q.
                                                               \tag{2}
\]
As $|U|\le n-3=2d+1$, this forces $q\le2$ and $|F|\le4$.

The involution $\epsilon$ centralizes $H$ and fixes $A$
pointwise, since $H$ contains pure $\Alt(A)$.
For $q=1$, it therefore fixes all of $U$.
For $q=2$, (2) forces $|U|=2d+1$. Let $E_0$ be the other
nontrivial orbit. If $\epsilon$ acted nontrivially on $E_0$,
it would be free there. Partition the $d-3$ paired colours
of $H\cap Q$ into those with an eligible edge in $E_0$
and the rest. The first set has at most $|E_0|/2-1$ colours:
project one eligible edge per colour to the $\epsilon$-pairs;
the actual $\eta$ belongs to every relevant ambient omission
and supplies the pair switches, so an alternate projected path
would contradict eligibility. The projected graph is a forest.
All eligible edges of the second set lie in the single orbit
$A\cap W$, so the complete selected-localization theorem bounds
that set by $(|A|-6)/2$. Their total is at most
\[
       \frac{|E_0|}{2}-1+\frac{|A|-6}{2}
        =\frac{2d+1-8}{2}=d-\frac72<d-3,
\]
a contradiction. Thus $\epsilon$ fixes $U$ in both cases.

The at most four points of $F$ include the fixed points $z,a$.
Oddness makes $\epsilon$ precisely one transposition on
the remaining points. Its pair $E$ is fixed pointwise by
$H$, and hence by $K$. Finally in either branch, $H$
centralizes $\epsilon$ and fixes $z,a$, so preserves $E$.
\end{proof}

\begin{lemma}[The actual penultimate mark is a double transposition]
Under the same hypotheses, $\sigma$ fixes all but four points.
\end{lemma}
\begin{proof}
Write $E=\{u,v'\}$ as in Lemma~1, and put $x=a^\sigma$.
As $P=\gen{H,\sigma}$ is transitive on $\Omega\setminus\{z\}$
and $H$ fixes $a$, we have $x\ne a$. Also $\sigma$ fixes $z$.
The group $K$ centralizes $\sigma$, so fixes pointwise
\[
                          z,a,E,E^\sigma,x.            \tag{3}
\]
It has at most six fixed points. Indeed its rank is $r-3$,
and its pure-prefix intersection with $Q$ removes five head
points while retaining its $q_K\ge1$ nontrivial orbits.
Every colour of $K\cap Q=G_{\{3,\ldots,r-4\}}$ is paired,
so the complete ambient interval inequality gives
\[
       2(r-6)\le |\operatorname{Supp}(K)|-5-q_K,
       \qquad |\operatorname{Supp}(K)|\ge2r-6.          \tag{4}
\]
Thus (3) cannot contain seven distinct points.

If $x\in E$, name it $u$. Then $u^\sigma=a$ and
$w=(v')^\sigma$ lies outside $\{z,a,u,v'\}$.
Injectivity and fixation of $z$ exclude the first three values;
$w=v'$ would make $\{a,u,v'\}$ invariant under $P$, since
$H$ fixes $a$ and preserves $E$. Thus the core of $\sigma$
is $(a\ u)(v'\ w)$ on five points, including the fixed $z$.

If $x\notin E$, the set $E^\sigma$ cannot equal $E$, since
$P$ would preserve $E$. If it were disjoint from $E$, then
$z,a,x,E,E^\sigma$ would be seven distinct fixed points of $K$,
contrary to (4). Their intersection therefore has one point.
As $\sigma$ is an involution, that point is fixed by it;
name it $v'$. Then $u^\sigma=w\notin\{z,a,x,u,v'\}$,
and the core is $(a\ x)(u\ w)$ on six points, including
the fixed $z,v'$.

In either case write $\sigma$ as the displayed two
transpositions times a disjoint involution $\nu$.
The latter is even, since $\sigma$ is even. The product
$\eta\sigma$ on the core has order five in the first case
and order three in the second; in the second case it consists
of two disjoint three-cycles. Therefore
$\nu=(\eta\sigma)^5$ or $(\eta\sigma)^3$, respectively.
It fixes $z$, so by (1) it belongs to $P$. Full IP gives
\[
                 \nu\in\gen{\eta,\sigma}\cap P=\gen\sigma.
\]
But $\nu$ fixes $a$ and $\sigma$ moves $a$. Thus $\nu=1$,
and $\sigma$ is literally the displayed double transposition.
\end{proof}

\begin{theorem}
The actual last three marks form the pure six-point
$A_5$ marking of source Table~2(3), read from the terminal end.
\end{theorem}
\begin{proof}
Apply the complete terminal-double-transposition theorem.
Its natural alternating prefix is exactly (1).
The group $H=G_{\le r-3}$ has its actual head orbit of
size at least four and contains its pure alternating group
by the acquired prefix-kernel theorem. Lemma~2 supplies
the literal double transposition $\rho_{r-2}$.
These are every hypothesis of that attachment theorem,
which gives the asserted three actual marks and removes
all exterior decoration through its full IP argument.
\end{proof}

\paragraph{Attribution and boundary.}
This supplies the terminal-singleton attachment in the route of
Proposition~5.32 of Cameron--Fernandes--Leemans--Mixer,
\emph{Highest rank of a polytope for $A_n$}, arXiv:1605.09173v1.
The complete odd-singleton theorem, pure-prefix kernels,
ambient localization and terminal attachment are acquired.
The present proof verifies the actual odd-tail embedding,
transfers its short fixed set past the supported left head,
and extracts the two literal final double transpositions.
It retains the full paired-interior hypothesis. Its conclusion
is an actual opposite head; excluding that branch and completing
all source routing are separate tasks. No whole-original,
question, numerical or formal credit is claimed.
\end{document}
